\documentclass[20pt]{beamer}
% \listfiles
 
\setbeamercolor{normal text}{bg=black, fg=white}
\setbeamercolor{frametitle}{fg=orange}
\setbeamertemplate{navigation symbols}{}
\setbeamertemplate{frametitle}[default][center]
\makeatletter 
\setbeamertemplate{frametitle continuation}{\gdef\beamer@frametitle{}}
\makeatother 
\setbeamerfont{frametitle}{size=\large}
\makeatletter
\define@key{beamerframe}{c}[true]{% centered
  \beamer@frametopskip=0pt plus 1fill\relax%
  \beamer@framebottomskip=0pt plus 1fill\relax%
  \beamer@frametopskipautobreak=0pt plus .4\paperheight\relax%
  \beamer@framebottomskipautobreak=0pt plus .6\paperheight\relax%
  \def\beamer@initfirstlineunskip{}%
}
\makeatother

\usepackage{fontspec}
\usepackage{polyglossia}
\setdefaultlanguage{malayalam}
% \setsansfont[Script=Malayalam, HyphenChar="0000]{Rachana}
\setsansfont[Script=Malayalam]{Rachana}
%\setsansfont[Script=Malayalam]{Meera}

% TODO: Justification ?
%\usepackage{ragged2e}
%\apptocmd{\frame}{}{\justifying}{} % Allow optional arguments after frame.
%\renewcommand{\raggedright}{\leftskip=0pt \rightskip=0pt plus 0cm} 

\setlength{\parskip}{0.5em}

\usepackage{graphicx}
\usepackage{tcolorbox}

\title{വിശുദ്ധ കുര്‍ബാന}
\subject{നോബുകാലം}
\author{github.com/SyroMalabar}

\newcommand{\Priest}[1]{\color{white}{#1}}
\newcommand{\People}[1]{\color{yellow}{#1}}
\newcommand{\Server}[1]{\color{cyan}{#1}}
\newcommand{\Notice}[1]{\color{lime}{(#1)}}

\newcommand{\SignOfCross}{പിതാവും പുത്രനും പരിശുദ്ധാത്മാവുമായ സര്‍വ്വേശ്വരാ, എന്നേക്കും.}
\newcommand{\Ammen}{\People{ആമ്മേന്‍.}}
\newcommand{\Peace}{\Server{സമാധാനം നമ്മോടുകൂടെ.}}
\newcommand{\PPeace}{\Priest{സമാധാനം + നിങ്ങളോടുകൂടെ.}}
\newcommand{\LetsPray}{\Server{നമുക്കൂ പ്രാര്‍ത്ഥിക്കാം,\\ സമാധാനം നമ്മോടുകൂടെ.}}
\newcommand{\ForgiveUs}{\People{കര്‍ത്താവേ, അങ്ങയുടെ  ദാസരുടെ പാപങ്ങളും അപരാധങ്ങളും ക്ഷമിക്കണമേ.}}
\newcommand{\DivyaKarunyaSthuthi}{\Server{ജീവിക്കുന്ന ദൈവത്തെ നിങ്ങള്‍ പ്രകീര്‍ത്തിക്കുവിന്‍.}}

\newcommand{\response}{കര്‍ത്താവേ}
\newcommand{\presponse}{\People{\response}}
\newcommand{\karosoosa}[1]{\Server{#1}\par\presponse\\\framebreak}

\newcommand{\AmbaramanavathamPPL}{
\People{
ദിവ്യത്മാവിന്‍ ഗീതികളാല്‍\\
ഹല്ലേലൂയ്യ ഗീതികളാല്‍\\
കര്‍ത്താവിന്‍ തിരുദിവസത്തിന്‍\\
നിര്‍മ്മലമാകുമനുസ്മരണം\\
കൊണ്ടാടാം ഇന്നീ വേദികയില്‍.}}

\begin{document}

\begin{frame}
\includegraphics[scale=0.15]{intro}
\end{frame}

\begin{frame}
\frametitle{കുരിശടയാളം}
\Priest{
പിതാവിന്റെയും പുത്രന്റെയും പരിശുദ്ധാത്മാവിന്റെയും നാമത്തില്‍.}

\Ammen
\end{frame}

\begin{frame}
\frametitle{മിശിഹായുടെ കല്പന}
\Priest{
അന്നാപ്പെസഹാത്തിരുനാളില്‍\\
കര്‍ത്താവരുളിയ കല്പനപോല്‍\\
തിരുനാമത്തില്‍ ചേര്‍ന്നീടാം\\
ഒരുമയോടീ ബലിയര്‍പ്പിക്കാം.}\par
\People{
അനുരഞ്ജിതരായ് തീര്‍ന്നീടാം\\
നവമൊരു പീഠമൊരുക്കീടാം\\
ഗുരുവിന്‍ സ്നേഹമോടീയാഗം\\
തിരുമുബാകെയണച്ചീടാം.}
\end{frame}

\begin{frame}
\frametitle{മാലാഖമാരുടെ കീര്‍ത്തനം}
\Priest{
അത്യുന്നതങ്ങളില്‍ ദൈവത്തിനു സ്തുതി.}\par
\Ammen\par
\Priest{
ഭൂമിയില്‍ മനുഷ്യർക്കു സമാധാനവും പ്രത്യാശയും എപ്പോഴും എന്നേക്കും.}\par
\Ammen
\end{frame}

\begin{frame}
\frametitle{മാലാഖമാരുടെ കീര്‍ത്തനം}
\Priest{
അത്യുന്നതമാം സ്വര്‍ലോകത്തില്‍ സര്‍വ്വേശനു സ്തുതിഗീതം.}\par
\People{
ഭൂമിയിലെങ്ങും മര്‍ത്യനു ശാന്തി പ്രത്യാശയുമെന്നേക്കും.}
\end{frame}

\begin{frame}[allowframebreaks]
\frametitle{സ്വര്‍ഗ്ഗസ്ഥനായ ഞങ്ങളുടെ പിതാവേ}
\Priest{സ്വർഗ്ഗസ്ഥനായ ഞങ്ങളുടെ പിതാവേ,}
\People{
അങ്ങയുടെ നാമം പൂജിതമാകണമേ./ അങ്ങയുടെ രാജ്യം വരണമേ./
അങ്ങയുടെ തിരുമനസ്സു സ്വര്‍ഗ്ഗത്തിലെ പോലെ ഭൂമിയിലുമാകണമേ./
ഞങ്ങള്‍ക്ക് ആവശ്യകമായ ആഹാരം/ ഇന്നു ഞങ്ങള്‍ക്കു തരണമേ.
ഞങ്ങളുടെ കടക്കാരോട് ഞങ്ങള്‍ ക്ഷമിച്ചിരിക്കുന്നത് പോലെ/
ഞങ്ങളുടെ കടങ്ങളും പാപങ്ങളും ഞങ്ങളോടും ക്ഷമിക്കണമേ./
ഞങ്ങളെ പ്രലോഭനത്തില്‍ ഉള്‍പ്പെടുത്തരുതേ/ ദുഷ്ടാരൂപിയില്‍ 
നിന്ന് ഞങ്ങളെ രക്ഷിക്കണമേ/ എന്തുകൊണ്ടെന്നാല്‍/ രാജ്യവും
ശക്തിയും മഹത്ത്വവും/ എന്നേക്കും അങ്ങയുടെതാകുന്നു.\par\Ammen}\\
\end{frame}

\begin{frame}
\People{
സ്വർഗ്ഗസ്ഥനായ ഞങ്ങളുടെ പിതാവേ,/ അങ്ങയുടെ മഹത്ത്വത്താല്‍/
സ്വര്‍ഗ്ഗവും ഭൂമിയും നിറഞ്ഞിരിക്കുന്നു./ മാലാഖമാരും മനുഷ്യരും/
അങ്ങ് പരിശുദ്ധന്‍, പരിശുദ്ധന്‍ പരിശുദ്ധന്‍ എന്ന് ഉദ്ഘോഷിക്കുന്നു.}
\end{frame}

\begin{frame}[allowframebreaks]
\frametitle{സ്വര്‍ഗ്ഗസ്ഥിതനാം താതാ നിന്‍}
\Priest{
സ്വര്‍ഗ്ഗസ്ഥിതനാം താതാ നിന്‍\\
നാമം പൂജിതമാകണമേ.\\
നിന്‍ രാജ്യം വന്നീടണമേ\\
പരിശുദ്ധന്‍ നീ പരിശുദ്ധന്‍.}\par
\People{
സ്വര്‍ഗ്ഗസ്ഥിതനാം താത നിന്‍\\
സ്തുതിതന്‍ നിസ്തുല മഹിമാവാല്‍\\
ഭൂസ്വര്‍ഗ്ഗങ്ങള്‍ നിറഞ്ഞു സദാ\\
പാവനമായി വിളങ്ങുന്നു.\par
\framebreak
വാനവമാനവവൃന്ദങ്ങള്‍\\
ഉദ്ഖോഷിപ്പു സാമോദം\\
പരിശുദ്ധന്‍ നീ എന്നെന്നും\\
പരിശുദ്ധന്‍ നീ പരിശുദ്ധന്‍.\par
സ്വര്‍ഗ്ഗസ്ഥിതനാം താത നിന്‍\\
നാമം പൂജിതമാകണമേ.\\ 
നിന്‍രാജ്യം വന്നീടണമേ\\
നിന്‍ ഹിതമിവിടെഭവിക്കണമേ.\par
\framebreak
സ്വര്‍ഗ്ഗത്തെന്നതുപോലുലകില്‍\\
നിന്‍ ചിത്തം നിറവേറണമേ.\\
ആവശ്യകമാമാഹാരം\\
ഞങ്ങള്‍ക്കിന്നരുളീടണമേ.\par
ഞങ്ങള്‍ കടങ്ങള്‍ പൊറുത്തതുപോല്‍\\
ഞങ്ങള്‍ക്കുള്ള കടം സകലം\\ 
പാപത്തിന്‍ കടബാദ്ധ്യതയും\\ 
അങ്ങു കനിഞ്ഞുപൊറുക്കണമേ.\par
\framebreak
ഞങ്ങള്‍ പരീക്ഷയിലോരുനാളും\\
ഉള്‍പ്പെടുവാനിടയാകരുതേ\\ 
ദുഷ്ടാരൂപിയില്‍ നിന്നെന്നും\\ 
ഞങ്ങളെ രക്ഷിച്ചരുളണമേ.\par
എന്തെന്നാലെന്നാളേക്കും\\ 
രാജ്യം ശക്തി മഹത്ത്വങ്ങള്‍\\
താവകമല്ലോ കര്‍ത്താവേ\\
ആമ്മേനാമ്മേനെന്നേക്കും.}
\end{frame}

\begin{frame}
\LetsPray
\end{frame}

% TODO: add prayers according to kalam
\begin{frame}[allowframebreaks]
\frametitle{പ്രാരംഭ പ്രാര്‍ത്ഥന-നോബുകാലം}
\Priest{ഞങ്ങളുടെ കര്‍ത്താവായ ദൈവമേ, അങ്ങയുടെ തിരുക്കുമാരന്റെ രക്ഷാകര 
രഹസ്യങ്ങളില്‍ ഞങ്ങളെ പങ്കുകാരാക്കുന്നതിന് ഞങ്ങള്‍ അങ്ങയെ 
സ്തുതിക്കുന്നു. പാപം മൂലം അകന്നുപോയ ഞങ്ങള്‍ അനുതപിക്കൂന്ന 
ഹൃദയതത്തോടെ പിതാവായ അങ്ങയെ സമീപിക്കുന്നു. മിശിഹായോടു 
ചേര്‍ന്നു എളിമയോടും  ഹൃദയവിശുദ്ധിയോടും കൂടി പരിഹാരത്തിന്റേയും 
ക്രിതജഞതയുടേയും ഈ ബലി അര്‍പ്പിക്കുവാന്‍ ഞങ്ങളെ യോഗ്യരാക്കണമേ.
സകലത്തിന്റേയും നാഥാ, എന്നെക്കും.}\par
\Ammen\\
\end{frame}


% Sundays and Ordinary Celebrations
\begin{frame}[allowframebreaks]
\frametitle{പ്രാരംഭ പ്രാര്‍ത്ഥന}
\Priest{
ഞങ്ങളുടെ കര്‍ത്താവായ ദൈവമേ, മനുഷ്യ വംശത്തിന്റെ
നവീകരണ ത്തിനും രക്ഷയ്ക്കും വേണ്ടി അങ്ങയുടെ പ്രിയപുത്രന്‍
കാരുണ്യപൂര്‍വ്വം നല്‍കിയ ദിവ്യരഹസ്യങ്ങളുടെ
പരികര്‍മ്മത്തിനു ബലഹീനരായ ഞങ്ങളെ ശക്തരാക്കണമേ.
സകലത്തിന്‍റെയും നാഥാ, എന്നേക്കും.} \Ammen
\end{frame}

% Other Celebrations
\begin{frame}[allowframebreaks]
\frametitle{പ്രാരംഭ പ്രാര്‍ത്ഥന}
\Priest{
ഞങ്ങളുടെ കര്‍ത്താവായ ദൈവമേ, അങ്ങയുടെ തിരുനാമത്തില്‍ 
ഉറച്ചുവിശ്വസിക്കുകയും ആ വിശ്വാസം പരമാര്‍ത്ഥതയോടെ 
ഏറ്റുപറയുകയും ചെയുന്നവരെ അങ്ങു ശക്തരാക്കണമേ.
ആത്മ ശരീരങ്ങളെ പവിത്രീകരിക്കുന്ന ഈ പരിഹാര രഹസ്യങ്ങള്‍ 
അവര്‍ വിശുദ്ധിയോടെ പരികര്‍മ്മം ചെയ്യട്ടെ. നിര്‍മ്മല
ഹൃദയത്തോടും വിശുദ്ധ വിചാരങ്ങളോടും കൂടെ അവര്‍
അങ്ങേക്കു പുരോഹിതശുശ്രൂഷ ചെയ്യുകയും അങ്ങു കനിഞ്ഞു 
നല്‍കിയ രക്ഷയെ പ്രതി നിരന്തരം അങ്ങയെ സ്തുതിക്കുകയും ചെയ്യട്ടെ.
\SignOfCross} \Ammen
\end{frame}

% Other Celebrations
\begin{frame}[allowframebreaks]
\frametitle{പ്രാരംഭ പ്രാര്‍ത്ഥന}
\Priest{
ഞങ്ങളുടെ കര്‍ത്താവായ ദൈവമേ, അങ്ങയുടെ മഹനീയ
ത്രിത്വത്തിന്റെ സംപൂജ്യമായ നാമത്തിനു സ്വര്‍ഗ്ഗത്തിലും ഭൂമിയിലും
എപ്പോഴും സ്തുതിയും ബഹുമാനവും കൃതജ്ഞതയും ആരാധനയും
ഉണ്ടായിരിക്കട്ടെ. \SignOfCross} \Ammen
\end{frame}

% TODO: add സങ്കീര്‍ത്തനം according to kalam
\begin{frame}[allowframebreaks]
\frametitle{സങ്കീര്‍ത്തനം-നോബുകാലം}

\Priest{കര്‍ത്താവേ, അങ്ങയുടെ ദാസര്‍ അങ്ങേയ്ക്കൂ നന്ദി പറയും. അങ്ങയുടെ നീതിമാന്മാര്‍ 
അങ്ങയെ വാഴ്ത്തുകയും ചെയ്യും.}\par
\People{ഞങ്ങളുടെ രക്ഷകനായ ദൈവമേ, ഞങ്ങളെ സഹായിക്കണമേ. അങ്ങയുടെ 
തിരുനാമത്തെപ്രതി ഞങ്ങളുടെ പാപങ്ങള്‍ ക്ഷമിക്കുകയും ചെയ്യണമെ.}\\\framebreak

\Priest{അങ്ങയുടെ രാജ്യത്തിന്റെ മഹത്വത്തെ അവര്‍ പ്രകീര്‍ത്തിക്കും.
അങ്ങയുടെ ശക്തിയെ അവര്‍ വര്‍ണിക്കും.}\par
\People{കര്‍ത്താവ് വീഴുന്നവരെയെല്ലം താങ്ങുന്നു. നിലം പതിച്ചവരെ എഴുന്നേല്‍പിക്കുന്നു.}\\\framebreak

\Priest{എല്ലാവരും കര്‍ത്താവില്‍ ദ്രിഷ്ട്ടി പതിച്ചിരിക്കുന്നു. എന്തെന്നാല്‍ യഥാസമയം അവര്‍ക്ക് 
അവിടുന്നു ആഹാരം നല്കുന്നു.}\par
\People{നിര്‍മ്മലഹൃദയത്തോടെ കര്‍ത്താവിനെ വിളിക്കൂന്നവര്‍ക്കു അവിടുന്നു സമീപസ്ഥനാകുന്നു. }\\\framebreak

\Priest{തന്റെ ഭക്തരുടെ ആഗ്രഹം അവിടുന്നു സഫലമാക്കുന്നു.}\par
\People{അവരുടെ പ്രാര്‍ത്ഥന കേട്ട് അവിടുന്നു അവരെ രക്ഷിക്കുന്നു.}\\~\\

\Priest{പിതാവിനും പുത്രനും പരിശുദ്ധാത്മാവിനും സ്തുതി.}\par
\People{ആതിമുതല്‍ എന്നേക്കും }\Ammen\\

\end{frame}


\begin{frame}[allowframebreaks]
\frametitle{സങ്കീര്‍ത്തനം}
\Priest{
കര്‍ത്താവേ, മമ രാജാവേ,\\
പാടും നിൻപുകളെന്നും ഞാൻ\\
സകലേശാ, നിൻ തിരുനാമം\\
വാഴ്ത്തീടും ഞാനനവരതം.}\par
\framebreak
\People{
കര്‍ത്താവേ, നിൻ സ്‌തുതി പാടും\\
അനുദിനമങ്ങയെ വാഴ്ത്തും ഞാൻ\\
നാഥൻ മഹിമ നിറഞ്ഞവനും\\
പാരം സ്തുത്യനുമെന്നെനും.\par
അനവദ്യൻ മമ നാഥാ, നിൻ\\
മഹിമയ്ക്കളവില്ലറിവൂ ഞാൻ\\
തലമുറ തലമുറയോടെന്നും\\
തവ കർമ്മങ്ങൾ വിവരിക്കും.}\\
\end{frame}

% not used
\iffalse
\begin{frame}[allowframebreaks]
\frametitle{ധൂപാശീര്‍വ്വാദം}
\Priest{
പിതാവും പുത്രനും പരിശുദ്ധാത്മാവുമായ സര്‍വേശ്വര,
അങ്ങയുടെ ബഹുമാനത്തി നായി ഞങള്‍ സമര്‍പ്പിക്കുന്ന ഈ ധുപം
അങ്ങയുടെ മഹനീയ ത്രിത്വത്തിന്‍റെ നാമത്തില്‍ + ആശീര്‍വ്വദിക്കപ്പെടട്ടെ.
ഇത് അങ്ങയുടെ പ്രസാദത്തിനും അങ്ങയുടെ അജഗണത്തിന്റെ
പാപങ്ങളുടെ മോചനത്തിനും കാരണമാകട്ടെ. എന്നേക്കും.} \Ammen
\end{frame}
\fi

\begin{frame}
\Peace
\end{frame}

% Sundays
\begin{frame}[allowframebreaks]
\frametitle{ഉത്ഥാനഗീതത്തിന് ഒരുക്കം}
\Priest{
ഞങ്ങളുടെ കര്‍ത്താവായ ദൈവമേ, അങ്ങയുടെ സ്നേഹത്തിന്റ
പരിമളം ഞങ്ങളില്‍ വീശുകയും അങ്ങയുടെ സത്യത്തിന്റെ
ജ്ഞാനം ഞങ്ങളുടെ ആത്മാക്കളെ പ്രകാശിപ്പിക്കുകയും
ചെയ്യുമ്പോള്‍ സ്വര്‍ഗ്ഗത്തില്‍നിന്നു പ്രത്യക്ഷനാകുന്ന
അങ്ങയുടെ തിരുക്കുമാരനെ സ്വീകരിക്കുവാന്‍ ഞങ്ങള്‍ക്ക്
ഇടയാകട്ടെ. സകല സൗഭാഗ്യങ്ങളും നന്മകളും നിറഞ്ഞ്
മുടിചൂടി നില്ക്കുന്ന സഭയില്‍ നിരന്തരം അങ്ങയെ സ്തുതിക്കുവാനും
മഹത്വപ്പെടുത്തുവാനും ഞങ്ങള്‍ യോഗ്യരാകട്ടെ.
എന്തുകൊണ്ടെന്നാല്‍, അങ്ങ് എല്ലാത്തിന്റെയും സ്രഷ്ടാവാകുന്നു.
സകലത്തിന്റെയും നാഥാ, എന്നേക്കും.} \Ammen
\end{frame}

% Ordinary Days
\begin{frame}[allowframebreaks]
\frametitle{ഉത്ഥാനഗീതത്തിന് ഒരുക്കം}
\Priest{
ഞങ്ങളുടെ കര്‍ത്താവായ ദൈവമേ, അങ്ങു നല്കിയുട്ടുള്ളതും
എന്നാല്‍ കൃത്യജ്ഞത പ്രകാശിപ്പിക്കുവാന്‍ ഞങ്ങള്‍ക്കു
കഴിയാത്തതുമായ എല്ലാ സഹായങ്ങള്‍ക്കും 
അനുഗ്രഹങ്ങള്‍ക്കും ആയി സകല സൗഭാഗ്യങ്ങളും
നന്മകളും നിറഞ്ഞ് മുടിചൂടി നില്‍ക്കുന്ന സഭയില്‍ ഞങ്ങള്‍
അങ്ങയെ നിരന്തരം സ്തുതിക്കുകയും മഹത്വപ്പെടുത്തുകയും ചെയ്യട്ടെ.
അങ്ങ് സകലത്തിന്റെയും നാഥനും സ്രഷ്ടാവുമാകുന്നു.
\SignOfCross} \Ammen
\end{frame}

\begin{frame}[allowframebreaks]
\frametitle{ഉത്ഥാനഗീതം}
\People{
സര്‍വാധിപനാം കര്‍ത്താവേ,\\
നിന്നെ വണങ്ങി നമിക്കുന്നു.\\
ഈശോനാഥാ, വിനയമോടെ\\
നിന്നെ നമിച്ചു പുകഴ്ത്തുന്നു.\par
മർത്യനു നിത്യമഹോന്നതമാം\\
ഉത്ഥാനം നീയരുളുന്നു.\\
അക്ഷയമാവനുടെ ആത്മാവി-\\
ന്നുത്തമരക്ഷയുമേകുന്നു.}
\end{frame}

\begin{frame}
\LetsPray
\end{frame}

\begin{frame}
\Priest{
എന്റെ കര്‍ത്താവേ, നീ സത്യമായും ഞങ്ങളുടെ ശരീരങ്ങളെ
ഉയര്‍പ്പിക്കുന്നവന്നും ആത്മാക്കളെ രക്ഷിക്കുന്നവനും ജീവനെ
നിത്യം പരിപാലിക്കുന്നവനുമാകുന്നു. ഞങ്ങള്‍ എപ്പോഴും
നിനക്ക് സ്തുതിയും കൃതജ്ഞതയും ആരാധനയും സമര്‍പ്പിക്കുവാന്‍
കടപ്പെട്ടവരാകുന്നു. സകലത്തിന്റെയും നാഥാ, എന്നേക്കും.}
\Ammen
\end{frame}

\begin{frame}
\frametitle{ത്രൈശുദ്ധ കീര്‍ത്തനം}
\People{
ശബ്ദമുയര്‍ത്തിപ്പാടിടുവിന്‍\\
സര്‍വ്വരൂമൊന്നായി പാടിടുവിന്‍\\
എന്നെന്നും ജീവിക്കും\\
സർവ്വേശ്വരനെ വാഴ്ത്തിടുവിന്‍.\par
പരിപാവനനാം സര്‍വ്വേശാ,\\
പരിപാവനനാം ബലവാനെ,\\
പരിപാവനനാം അമര്‍ത്യനെ,\\
നിന്‍കൃപ ഞങ്ങള്‍ക്കേകണമേ.}
\end{frame}

\begin{frame}
\frametitle{വചനശുശ്രൂഷ}
\Priest{
വിശുദ്ധരിൽ സംപ്രീതനായി വസിക്കുന്ന പരിശുദ്ധനും
സ്തുത്യര്‍ഹനും ബലവാനും അമര്‍ത്യനുമായ കര്‍ത്താവേ,
അങ്ങയുടെ സ്വഭാവത്തിനൊത്തവിധം എപ്പോഴും
ഞങ്ങളെ കടാക്ഷിക്കുകയും അനുഗ്രഹിക്കുകയും ഞങ്ങളോടു
കരുണ കാണിക്കുകയും ചെയ്യണമേ. \SignOfCross} \Ammen
\end{frame}

\begin{frame}
\frametitle{പഴയനിയമ വായന}
\Server{
സഹോദരരേ, നിങ്ങള്‍ ഇരുന്നു ശ്രദ്ധയോടെ ദൈവവചനം കേള്‍ക്കുവിന്‍.}\par
\Notice{സമൂഹം ഇരിക്കുന്നു.}\par
\People{ദൈവമായ കര്‍ത്താവിനു സ്തുതി.}
\end{frame}

\begin{frame}[allowframebreaks]
\frametitle{സ്തുതിഗീതം}
\Priest{
അംബരമനവരതം\\
ദൈവമഹത്വത്തെ\\
വാഴ്ത്തിപ്പാടുന്നു.}\par
\AmbaramanavathamPPL\par
\framebreak
\Priest{
തൻ കരവിരുതല്ലോ\\
വാനവിതാനങ്ങൾ\\
ഉദ്ഘോഷിക്കുന്നു.}\par
\AmbaramanavathamPPL\par
\framebreak
\Priest{
പകലുകള്‍ പകലുകളോ-\\
ടവിരതമവിടുത്തെ\\
പുകളുര ചെയ്യുന്നു.}\par
\AmbaramanavathamPPL\par
\framebreak
\Priest{
നിത്യപിതാവിനും\\
സുതനും റൂഹായ്ക്കും\\
സ്തുതിയുണ്ടാകട്ടെ.}\par
\AmbaramanavathamPPL\par
\framebreak
\Priest{
ആദിയിലെപ്പോലെ\\
ഇപ്പോഴുമെപ്പോഴും\\
എന്നേക്കും ആമ്മേന്‍.}\par
\AmbaramanavathamPPL
\end{frame}

\begin{frame}
\Server{ഹല്ലേലൂയ്യ ഹല്ലേലൂയ്യ ഹല്ലേലൂയ്യ;}\par
\LetsPray
\end{frame}

% Sundays
\begin{frame}[allowframebreaks]
\Priest{
 ഞങ്ങളുടെ കര്‍ത്താവായ ദൈവമേ, അങ്ങയുടെ
 ജീവദായകവും ദൈവികവുമായ കല്പനകളുടെ മധുരസ്വരം
 ശ്രവിക്കുന്നതിനും ഗ്രഹിക്കുന്നതിനും ഞങ്ങളുടെ ബുദ്ധിയെ
 പ്രകാശിപ്പിക്കണമേ. അതുവഴി ആത്മശരീരങ്ങള്‍ക്ക്
ഉപകരിക്കുന്ന സ്നേഹവും ശരണവും രക്ഷയും ഞങ്ങളില്‍
ഫലമണിയുന്നതിന്നും നിരന്തരം ഞങ്ങള്‍ അങ്ങയെ
സ്തുതിക്കുന്നതിനും അങ്ങയുടെ കാരുണ്യത്താലും
അനുഗ്രഹത്താലും ഞങ്ങളെ സഹായിക്കണമേ.
\SignOfCross} \Ammen
\end{frame}

% TODO
% Ordinary Days

\begin{frame}
\frametitle{ലേഖന വായന}
\People{നമ്മുടെ കര്‍ത്താവായ മിശിഹായ്ക്കു സ്തുതി.}\par
\Notice{സമൂഹം എഴുന്നേല്‍ക്കുന്നു.}
\end{frame}

\begin{frame}[allowframebreaks]
\frametitle{ഹല്ലേലൂയ്യ ഗീതം}
\People{
ഹല്ലേലൂയ്യ പാടീടുന്നേന്‍\\
ഹല്ലേലൂയ്യ ഹല്ലേലൂയ്യ.\par
നല്ലോരാശയമെന്‍ മനതാരില്‍\\
വന്നു നിറഞ്ഞു തുളുബീടുന്നു.\par
രാജാവിന്‍ തിരുമുമ്പില്‍ കീര്‍ത്തന\\
മധുവായ് ഞാനതൊഴുക്കീടട്ടെ.\\
\framebreak
ഏറ്റമനുഗ്രഹപൂരിതനാം കവി തന്‍ \\
തൂലികപോലെന്‍നാവിപ്പോള്‍.\par
താതനുമതുപോല്‍ സുതനും\\
പരിശുദ്ധാത്മാവിനും സ്തുതിയു യരട്ടെ.\par
ആദിമുതല്ക്കേയിന്നും നിത്യവു-\\
മായി ഭവിചീടട്ടെ, ആമ്മേന്‍.\par
ഹല്ലേലൂയ്യ പാടീടുന്നേന്‍\\
ഹല്ലേലൂയ്യ ഹല്ലേലൂയ്യ.}
\end{frame}

% not used
\iffalse
\begin{frame}[allowframebreaks]
\frametitle{സുവിശേഷവായനയ്ക്ക് ഒരുക്കം}
\Priest{
ലോകത്തിന്‍റെ പ്രകാശവും സകലത്തിന്‍റെയും ജീവനുമായ 
മിശിഹായെ, നിന്നെ ഞങ്ങളുടെ പക്കലേക്കയച്ച 
അനന്തകാരുണ്യത്തിന് എന്നേക്കും സ്തുതി. ആമ്മേന്‍.}
\end{frame}
\fi

\begin{frame}
\frametitle{സുവിശേഷ വായന}
\Server{നമുക്കു ശ്രദ്ധാപൂര്‍വ്വം നിന്ന് പരിശുദ്ധ സുവിശേഷം ശ്രവിക്കാം.}\\~\\
\PPeace\\
\People{അങ്ങയോടും അങ്ങയുടെ ആത്മാവോടും കൂടെ.}\\~\\
\People{നമ്മുടെ കര്‍ത്താവായ മിശിഹായ്ക്കു സ്തുതി.}
\end{frame}

% TODO കാറോസൂസ

\begin{frame}[allowframebreaks]
\frametitle{കാറോസൂസ-നോബുകാലം}
\renewcommand{\response}{കര്‍ത്താവേ, ഞങ്ങളുടെ പ്രാര്‍ത്ഥന കേള്‍ക്കണമേ.}

\karosoosa{നമുക്കെല്ലാവര്‍ക്കും അനുതാപത്തോടും ഭക്തിയോടും കൂടി നിന്ന് \response എന്ന് അപേക്ഷിക്കാം.}

\karosoosa{പരിശുദ്ധാത്മാവിനാല്‍ വനത്തിലേയ്ക്കു നയിക്കപ്പെടുകയും നാല്പതു രാവും നാലപതു 
പകലും തപസനുഷുിക്കുകയും ചെയ്ത കര്‍ത്താവേ, തപസ്സിന്റെയും പ്രായഷിച്ചിത്തത്തിന്റെയും അരൂപിയാല്‍
ഞങ്ങളെ നിറയ്ക്കണമെന്ന് നിന്നൊട് ഞങ്ങള്‍ പ്രാര്‍ത്ഥിക്കൂന്നു.}

\karosoosa{ലോകത്തിന്റെ പാപങ്ങള്‍ നീക്കുന്ന ദിവ്യകുഞ്ഞാടായ മിശിഹായേ, ഞങ്ങളുടെയും 
മറ്റുള്ളവരുടെയും പാപങ്ങള്‍ക്കൂ പരിഹാരമനുഷുിക്കുവാന്‍ ഞങ്ങളെ ശക്തരാക്കണമെന്ന് ഞങ്ങള്‍ പ്രാര്‍ത്ഥിക്കൂന്നു.}

\karosoosa{സ്നേഹത്തിന്റെ ബലിവസ്തുവായ മിശിഹായേ, ഞങ്ങളുടെ ശരീരത്തെയും അതിന്റെ 
ദുര്‍വാസനകളെയും നിന്നോടുകൂടെ ക്രൂശിക്കൂന്നതിന് ഞങ്ങളെ അനുഗ്രഹിക്കണമെന്ന് ഞങ്ങള്‍ പ്രാര്‍ത്ഥിക്കൂന്നു.}

\karosoosa{സാത്താനെ പരാജയപ്പടുത്തി ഞങ്ങള്‍ക്കൂ മാതൃകയും മനോധൈര്യവും നല്കിയ മിശിഹായേ, 
പ്രലോഭനങ്ങളെ അതിജീവിക്കൂവാന്‍ ഞങ്ങളെ സക്തരാകണമെന്ന് ഞങ്ങള്‍ പ്രാര്‍ത്ഥിക്കൂന്നു.}

\karosoosa{പ്രാര്‍ത്ഥനയുടെയും പരിത്യാഗത്തിന്റെയും മാര്‍ഗ്ഗത്തില്‍ ഞങ്ങളെ നയിക്കുവാന്‍ 
ഞങ്ങളുടെ പരിശുദ്ധ പിതാവ് മാര്‍ ... പാപ്പായേയും ഞങ്ങളുടെ സഭയുടെ തലവനും പിതാവുമായ
മാര്‍ ... ശ്രേഷ്ഠമെത്രാപ്പോലീത്തായേയും ഞങ്ങളുടെ മേലദ്ധ്യക്ഷനും പിതാവുമായ മാര്‍ ... 
മെത്രാനേയും മറ്റ് അജപാലകരേയും സമ്രിദ്ധമായി അനുഗ്രഹിക്കണമെന്ന് ഞങ്ങള്‍ പ്രാര്‍ത്ഥിക്കൂന്നു.}

\renewcommand{\response}{ഞങ്ങളുടെ ദൈവമായ കര്‍ത്താവേ, അങ്ങേക്കു ഞങ്ങള്‍ സമര്‍പ്പിക്കുന്നു.}
\karosoosa{നമുക്കെലാവര്‍ക്കും നമ്മെയും നമോരുരത്തരെയും പിതാവിനും പുത്രനും പരിശുദ്ധാത്മാവിനും സമര്‍പ്പിക്കാം.}

\end{frame}


\begin{frame}
%\frametitle{കാഴ്ചസമര്‍പ്പണഗീതം} % -> cant fit on one slide
\Server{കര്‍ത്താവില്‍ ഞാന്‍ ദൃഡമായി ശരണപ്പെട്ടു}\\
\People{
മിശിഹാ കര്‍ത്താവിന്‍\\
തിരുമെയ്‌ നിണവുമിതാ\\
പാവന ബലിപീഠേ\\
സ്നേഹഭയങ്ങളോടണയുക നാ-\\
മഖിലരുമോന്നായ് സന്നിധിയില്‍\\
വാനവനിരയോടു ചേര്‍ന്നേവം \\
പാടാം ദൈവം പരിശുദ്ധന്‍\\
പരിശുദ്ധന്‍, നിത്യം പരിശുദ്ധന്‍}
\end{frame}

\begin{frame}[allowframebreaks]
\Priest{
താതനുമതുപോലത്മജനും ദിവ്യറുഹായ്ക്കും സ്തുതിയെന്നും
ദൈവാംബികയാകും കന്യാമറിയത്തെ 
സാദരമോർത്തീടാം പാവനമീബലിയില്‍}\par
\People{
ആദിയിലെപ്പോല്‍ എന്നെന്നേക്കും
ആമ്മേനാമ്മേന്‍.\\
സുതനുടെ പ്രേഷിതരെ, ഏകജസ്നേഹിതരെ,
ശാന്തി ലഭിച്ചിടുവാന്‍ നിങ്ങള്‍ പ്രാര്‍ത്ഥിപ്പിന്‍.}\\
\framebreak
\Priest{
സര്‍വ്വരുമൊന്നായ്‌ പാടീടട്ടെ 
ആമ്മേനാമ്മേന്‍.
മാർത്തോമ്മായേയും, നിണസാക്ഷികളെയും 
സത്ക്കർമ്മികളേയും, ബലിയിതിലോര്‍ത്തീടാം.}\par
\People{
നമ്മുടെകൂടെ ബലവാനാം
കര്‍ത്താവെന്നെന്നേയ്ക്കും.
രാജാവാം ദൈവം, നമ്മോടൊത്തെന്നും 
യാക്കോബിന്‍ ദൈവം, നമ്മുടെ തുണയെന്നും.}\\
\framebreak
\Priest{
ചെറിയവരെല്ലാം വലിയവരൊപ്പം
കാത്തുവസിക്കുന്നു.\\
മൃതരെല്ലാവരും നിൻ, മഹിതോത്ഥാനത്തിൽ\\
ശരണം തേടുന്നു ഉത്ഥിതരായിടുവാന്‍.}\par
\People{
തിരുസ്സന്നിധിയില്‍  ഹൃദയഗതങ്ങള്‍
ചൊരിയുവിനെന്നേക്കും\\
നോമ്പും പ്രാര്‍ത്ഥനയും, പശ്ചാത്താപവുമായ്
ത്രിത്വത്തെ മോദാല്‍, നിത്യം വാഴ്ത്തീടാം.}\\
\end{frame}

\begin{frame}[allowframebreaks]
\frametitle{വിശ്വാസപ്രഖ്യാപനം}
\Priest{
സര്‍വ്വശക്തനും പിതാവുമായ ഏക ദൈവത്തില്‍ ഞങ്ങള്‍ വിശ്വസിക്കുന്നു.}\par
\People{
ദൃശ്യവും അദൃശ്യവുമായ/ സകലത്തിന്റെയും സ്രഷ്ടാവില്‍/
ഞങ്ങള്‍ വിശ്വസിക്കുന്നു./ ദൈവത്തിന്റെ ഏകപുത്രനും/
സകല സൃഷ്ടികള്‍ക്കും മുമ്പുള്ള ആദ്യജാതനും/
യുഗങ്ങള്‍ക്കെല്ലാം മുമ്പ് പിതാവില്‍ \\ \framebreak 
നിന്ന് ജനിച്ചവനും/ എന്നാല്‍ സൃഷ്ട്ടിക്കപെടാത്തവനും/
എകകര്‍ത്താവുമായ ഈശോ/ മിശിഹായില്‍ ഞങ്ങള്‍ വിശ്വസിക്കുന്നു./
അവിടുന്നു സത്യദൈവത്തില്‍ നിന്നുള്ള സത്യദൈവവും/
പിതാവിനോടു കൂടെ ഏകസത്തയുമാകുന്നു./
അവിടുന്നു വഴി പ്രപഞ്ചം സംവിധാനം ചെയ്യപ്പെടുകയും/
എല്ലാം സൃഷ്ടിക്കപ്പെടുകയും ചെയ്തു./ മനുഷ്യരായ നമുക്ക് വേണ്ടിയും/
\\ \framebreak 
നമ്മുടെ രക്ഷയ്ക്കുവേണ്ടിയും/ അവിടുന്നു 
സ്വര്‍ഗ്ഗത്തില്‍ നിന്നിറങ്ങി/ പരിശുദ്ധാത്മാവിനാല്‍/
കന്യകാമറിയത്തില്‍ നിന്ന് ശരീരം സ്വീകരിച്ച്/മനുഷ്യനായിപ്പിറന്നു.
പന്തിയോസ് പീലാത്തോസിന്‍റെ കാലത്ത്
പീഡകള്‍ സഹിക്കുകയും/ സ്ലീവായില്‍ തറയ്ക്കപ്പെട്ടു 
മരിക്കുകയും/ സംസ്കരിക്കപ്പെടുകയും/
എഴുതപ്പെട്ടിരിക്കുന്നതു പോലെ/\\ \framebreak
മൂന്നാം ദിവസം ഉയിര്‍ത്തെഴുന്നേല്ക്കുകയും ചെയ്തു./
അവിടുന്നു സ്വര്‍ഗ്ഗത്തിലേക്ക് എഴുന്നള്ളി/
പിതാവിന്റെ വലതു ഭാഗത്തിരി ക്കുന്നു./
മരിച്ചവരെയും ജീവിക്കുന്നവരെയും വിധിക്കുവാന്‍/
അവിടുന്നു വീണ്ടും വരുവാനിരിക്കുന്നു./
പിതാവില്‍നിന്നും പുത്രനില്‍നിന്നും പുറപ്പെടുന്ന/
സത്യാത്മാവും ജീവദാതാവുമായ/\\ \framebreak
ഏക പരിശുദ്ധാത്മാവിലും ഞങ്ങള്‍ വിശ്വസിക്കുന്നു.
ഏകവും പരിശുദ്ധവും ശ്ലൈഹീകവും സാര്‍വ്വത്രികവുമായ/
സഭയിലും ഞങ്ങള്‍ വിശ്വസിക്കുന്നു./
പാപമോചനത്തിനായുള്ള ഏക മാമ്മോദീസായും/
ശരീരത്തിന്റെ ഉയിര്‍പ്പും/ നിത്യായുസ്സും ഞങ്ങള്‍ ഏറ്റു
പറയുകയും ചെയ്യുന്നു.\par \Ammen}
\end{frame}


\begin{frame}[allowframebreaks]
\frametitle{അനുസ്മരണ പ്രാര്‍ത്ഥന}
\Priest{
സകലത്തിന്റെയും നാഥനായ ദൈവം തന്റെ സ്തുതികൾ ആലപിക്കുവാൻ നിന്നെ + ശക്തനാക്കട്ടെ.}\par
\Server{
നമുക്ക് പ്രാര്‍ത്ഥിക്കാം, സമാധാനം നമ്മോടുകൂടെ.
........... 
സ്വര്‍ഗ്ഗരാജ്യത്തില്‍ നിത്യജീവനും കാരണമാകട്ടെ.}
\end{frame}

\begin{frame}[allowframebreaks]
\frametitle{മദ്ബഹാപ്രവേശനം}
\Priest{
കര്‍ത്താവേ ആശീര്‍വ്വദിക്കണമേ. എന്റെ സഹോദരരെ,
ഈ കുര്‍ബാന എന്റെ കരങ്ങള്‍വഴി പൂര്‍ത്തിയാകുവാന്‍
നിങ്ങള്‍ എനിക്കുവേണ്ടി പ്രാര്‍ത്ഥിക്കുവിന്‍.}\par
\People{
സകലത്തിന്റെയും നാഥനായ ദൈവം/ തന്റെ അഭീഷ്ടം നിറവേറ്റുവാന്‍/
നമ്മെ ശക്തരാക്കട്ടെ./ ഈ കുര്‍ബാന അവിടുന്ന് സ്വീകരിക്കട്ടെ/
അങ്ങേക്കും ഞങ്ങള്‍ക്കും/ ലോകം മുഴുവനും വേണ്ടി/
അങ്ങു സമര്‍പ്പിക്കുന്ന ഈ ബലിയില്‍/ അവിടുന്ന് 
സംപ്രീതനാവുകയും ചെയ്യട്ടെ.\par\Ammen}
\end{frame}

\begin{frame}[allowframebreaks]
\frametitle{ഒന്നാം പ്രമാണജപം}
\Priest{
ഞങ്ങളുടെ കര്‍ത്താവായ ദൈവമേ, അങ്ങു ഞങ്ങളുടെമേല്‍
വര്‍ഷിച്ച സമൃദ്ധമായ അനുഗ്രഹങ്ങളെ പ്രതി ഞങ്ങള്‍ അങ്ങേക്കു 
നന്ദി പറയുന്നു. അങ്ങയുടെ അഭിഷിക്തന്റെ ശരീര രക്തങ്ങളാകുന്ന
ദിവ്യരഹസ്യങ്ങളുടെ ശുശ്രൂഷകരാകുവാന്‍ എളിയവരും
പാപികളുമായിരുന്നിട്ടും ഞങ്ങളെ അങ്ങു കാരുണ്യാതിരേകത്താല്‍
യോഗ്യരാക്കി. അങ്ങു നല്കിയ ഈ ദാനം തികഞ്ഞ സ്നേഹത്തോടും
ഉറച്ച വിശ്വാസത്തോടും കു‌ടെ പരികര്‍മ്മം ചെയ്യുവാന്‍ ഞങ്ങളെ
ശക്തരാക്കണമെന്ന് അങ്ങയോടു ഞങ്ങള്‍ അപേക്ഷിക്കുന്നു.\par
ഞങ്ങള്‍ അങ്ങേക്കു സ്തുതിയും ബഹുമാനവും കൃതജ്ഞതയും ആരാധനയും 
സമര്‍പ്പിക്കുന്നു. ഇപ്പോഴും എപ്പോഴും എന്നേക്കും. \Ammen}\\
\end{frame}

\begin{frame}
\frametitle{സമാധാനാശംസ}
\PPeace \par
\People{അങ്ങയോടും അങ്ങയുടെ ആത്മാവോടും കൂടെ.} \par
\Server{സഹോദരരേ, മിശിഹായുടെ സ്നേഹത്തില്‍ നിങ്ങള്‍ സമാധാനം ആശംസിക്കുവിന്‍.}
\end{frame}

\begin{frame}[allowframebreaks]
\Server{
നമുക്കെല്ലാവര്‍ക്കും അനുതാപത്തോടും വിശുദ്ധിയോടുംകൂടെ കര്‍ത്താവിനു
നന്ദി പറയുകയും അവിടുത്തെ വിളിച്ചപേക്ഷിക്കുകയും ചെയ്യാം.
നിങ്ങള്‍ ആദരപൂര്‍വ്വം ഇവിടെ അനുഷ്ഠിക്കപ്പെടുന്നവയെ സൂക്ഷിച്ചു 
വീക്ഷിക്കുവിന്‍. ഭയഭക്തി ജനകമായ രഹസ്യങ്ങള്‍ കൂദാശ ചെയ്യപ്പെടുന്നു.
പുരോഹിതന്‍ തന്റെ മാദ്ധ്യസ്ഥ്യംവഴി സമാധാനം സമൃദ്ധമാകുന്നതിനു
വേണ്ടി പ്രാര്‍ത്ഥിക്കുവാന്‍ തുടങ്ങുന്നു. കണ്ണുകള്‍ താഴ്ത്തി, വിചാരങ്ങള്‍ 
സ്വര്‍ഗ്ഗത്തില്ലേക്ക് ഉയര്‍ത്തി, നിശബ്ദരായി ഏകാഗ്രതയോടും
ഭക്തിയോടും കൂടെ നിങ്ങള്‍ ഹൃദയത്തില്‍ പ്രാര്‍ത്ഥിക്കുവിന്‍. \Peace}
\end{frame}

\begin{frame}[allowframebreaks]
\frametitle{ആശംസാപ്രാര്‍ത്ഥന}
\Priest{
മിശിഹാ കര്‍ത്താവിന്‍ കൃപയും\\
ദൈവപിതാവിന്‍ സ്നേഹമതും\\
റൂഹാതന്‍ സഹവാസവുമീ +\\
നമ്മോടൊത്തുണ്ടാകട്ടെ.}\par
\Ammen\\ \framebreak
\Priest{
ഉയരങ്ങളിലേക്കുയരട്ടെ\\
ഹൃദയവികാരവിചാരങ്ങള്‍\\
ഉയരങ്ങളിലേക്കുയരട്ടെ\\
ഹൃദയവികാരവിചാരങ്ങള്‍.}\par
\People{
പൂര്‍വ്വപിതാവാമബ്രാഹം,\\
ഇസഹാക്ക് യാക്കോബെന്നിവര്‍തന്‍\\
ദൈവമേ, നിത്യം ആരാധ്യന്‍ \\
രാജാവേ, നിന്‍ സന്നിധിയില്‍}\\ \framebreak
\Priest{
അഖിലചരാചരകർത്താവാം\\
ദൈവത്തിന്നു ബലിയർപ്പിപ്പൂ.}\par
\People{
ന്യായവുമാണതു യുക്തവുമാം.\\
ന്യായവുമാണതു യുക്തവുമാം.}\\[2\baselineskip]
\Peace
\end{frame}

\begin{frame}[allowframebreaks]
\frametitle{രണ്ടാം പ്രണാമജപം}
\Priest{
പിതാവും പുത്രനും പരിശുദ്ധാത്മാവുമായ സര്‍വ്വേശ്വരാ,
അങ്ങയുടെ മഹത്ത്വമേറിയ ത്രിത്വത്തിന്റെ ആരാധ്യമായ
നാമം എല്ലാ അധരങ്ങളിലും നിന്ന് സ്തുതിയും എല്ലാ
നാവുകളിലുംനിന്ന് കൃതജ്ഞതയും എല്ലാ സൃഷ്ടികളില്‍
നിന്ന് പുകഴ്ചയും അര്‍ഹിക്കുന്നു. എന്തുകൊണ്ടെന്നാല്‍,
അങ്ങ് ലോകത്തേയും അതിലുള്ള സകലത്തെയും കനിവോടെ
സൃഷ്ട്ടിക്കുകയും മനുഷ്യവംശത്തോട് അളവറ്റ കൃപ കാണിക്കുകയും
ചെയ്തു. സ്വര്‍ഗ്ഗവാസികളുടെ ആയിരങ്ങളും മാലഖാമാരുടെ
പതിനായിരങ്ങളും മഹോന്നതനായ അങ്ങയെ കുമ്പിട്ടാരാധിക്കുന്നു.
\\ \framebreak
അഗ്നിമയന്മാരും അരൂപികളുമായ സ്വര്‍ഗ്ഗീയ സൈന്യങ്ങല്‍
അങ്ങയുടെ നാമം പ്രകീര്‍ത്തിക്കുന്നു. പരിശുദ്ധരും അരൂപികളുമായ
ക്രോവേന്മാരോടും സ്രാപ്പേന്മാരോടും ചേര്‍ന്ന് നാഥനായ അങ്ങേക്ക്
അവര്‍ ആരാധന സമര്‍പ്പിക്കുന്നു.}
\end{frame}

\begin{frame}[allowframebreaks]
\frametitle{ഓശാനാഗീതം}
\Priest{
ഒന്നായ് ഉച്ചസ്വരത്തിലവര്‍\\
തിരുസന്നിധിയില്‍ അനവരതം\\ 
സ്തുതിഗീതങ്ങള്‍ പാടുന്നു.}\par
\People{
ദൈവം നിത്യമഹത്വത്തിന്‍\\
കര്‍ത്താവെന്നും പരിശുദ്ധന്‍\\
ബലവാനീശന്‍ പരിശുദ്ധന്‍\\
മണ്ണും വിണ്ണും നിറയുന്നു.\\ \framebreak
മന്നവനവനുടെ മഹിമാകളാല്‍\\
ഉന്നതവീഥിയിലോശാന\\
ദാവീദിന്‍ സുതനോശാന\\
കര്‍ത്താവിന്‍ തിരുനാമത്തില്‍\\
വന്നവനും യുഗരാജാവായ്\\
വീണ്ടും വരുവോനും ധന്യന്‍\\
ഉന്നതവീഥിയിലോശാന.}\\~\\
\Notice{സമൂഹം മുട്ടുകുത്തുന്നു}
\end{frame}

% OR
\begin{frame}[allowframebreaks]
\frametitle{ഓശാനാഗീതം}
\People{
ബലവാനായ ദൈവമേ\\
പരിശുദ്ധന്‍.. പരിശുദ്ധന്‍..\\
കര്‍ത്താവായ ദൈവമേ\\
പരിശുദ്ധന്‍.. പരിശുദ്ധന്‍..\\
വിണ്ടല നാഥാ ഓശാനാ ....\\
മഹിതല നാഥാ ഓശാനാ ....\\
സര്‍വ്വചരാചര സൃഷ്ടാവേ പരിപാലകനെ വാഴ്ത്തിടാം
\hfill പരിശുദ്ധന്‍ ....}
\end{frame}

\begin{frame}[allowframebreaks]
\frametitle{മൂന്നാം പ്രണാമജപം}
\Priest{
കര്‍ത്താവായ ദൈവമേ, സ്വര്‍ഗ്ഗീയ ഗണങ്ങളോടു കൂടെ
അങ്ങേക്കു ഞങ്ങള്‍ കൃതജ്ഞത സമര്‍പ്പിക്കുന്നു.
അങ്ങയില്‍ മറഞ്ഞിരിക്കുന്ന അത്മജാതനും അങ്ങയോട്
സദൃശനും അങ്ങയില്‍ നിന്നുള്ള പ്രകാശവും അങ്ങയുടെ
സത്തയുടെ പ്രതിച്ഛായയുമായ വചനം ആകുന്ന ദൈവത്തെ
ഞങ്ങള്‍ വാഴ്ത്തുന്നു. അങ്ങയോടുള്ള സമാനത നില നിര്‍ത്തേണ്ട
കാര്യമായി പരിഗണിക്കാതെ അവിടുന്നു തന്നെത്തന്നെ
ശൂന്യനാക്കി, ദാസന്റെ രൂപം സ്വീകരിച്ച് വിവേകവും ബുദ്ധിയും ഉള്ള
അമര്‍ത്യമായ ആത്മാവോടും മര്‍ത്യമായ ശരീരത്തോടുംകൂടെ
പരിപൂര്‍ണ്ണ മനുഷ്യനായി സ്ത്രീയില്‍നിന്നു ജാതനായി.
നിയമത്തിന് അധീനരായവരെ ഉദ്ധരിക്കുവാന്‍ \\ \framebreak
നിയമത്തിന്നു വിധേയനാവുകയും ഞങ്ങളുടെ രക്ഷയുടെ സ്മാരകം
ഏര്‍പ്പെടുത്തുകയും ചെയ്തു. രക്ഷാകരമായ ഈ  രഹസ്യം
തിരുസന്നിധിയില്‍ ഞങ്ങള്‍ അര്‍പ്പിക്കുന്നു.

കര്‍ത്താവായ ദൈവമേ, അങ്ങയുടെ പ്രിയപുത്രന്‍ ഞങ്ങളെ
പഠിപ്പിച്ചതു പോലെ അവിടുത്തെ പീഡാനുഭവ ത്തിന്റെ സ്മരണ
ഞങ്ങള്‍ ആചരിക്കുന്നു. താന്‍ ഏല്പിച്ചു കൊടുക്കപ്പെട്ട രാത്രിയില്‍  \\ \framebreak
ഈശോ നിര്‍മ്മലമായ തൃക്കരങ്ങളില്‍ അപ്പമെടുത്തു സ്വര്‍ഗ്ഗത്തിലേക്ക്, 
ആരാധ്യനായ പിതാവേ, അങ്ങയുടെ പക്കലേക്ക് കണ്ണുകള്‍ ഉയര്‍ത്തി
വാഴ്ത്തി, വിഭജിച്ച്, ശിഷ്യന്മാര്‍ക്കു നല്‍കി കൊണ്ട് അരുളിച്ചെയ്തു.
ഇത് പാപ മോചനത്തിനായി നിങ്ങള്‍ക്ക് വേണ്ടി വിഭജിക്കപ്പെടുന്ന
എന്റെ ശരീരമാകുന്നു. നിങ്ങളെല്ലാവരും ഇതില്‍ നിന്ന് വാങ്ങി ഭക്ഷിക്കുവിന്‍.}
\Ammen
\end{frame}

\begin{frame}
\Priest{
അപ്രകാരംതന്നെ കാസയുമെടുത്ത് കൃതജ്ഞതാ സ്തോത്രം ചെയ്ത്,
വാഴ്ത്തി + അവര്‍ക്കു നല്കികൊണ്ട് അരുളിച്ചെയ്തു. ഇത് പാപമോചനത്തിനായ്
അനേകര്‍ക്ക്‌ വേണ്ടി ചിന്തപ്പെടുന്ന, പുതിയ ഉടമ്പടിയിലെ എന്റെ
രക്ത്തമാകുന്നു. നിങ്ങളെല്ലാവരും ഇതില്‍നിന്നു വാങ്ങി പാനം ചെയ്യുവിന്‍.}\par
\Ammen
\end{frame}

\begin{frame}
\Priest{
ഞാന്‍ ഈ ചെയ്തത് നിങ്ങള്‍ എന്റെ നാമത്തില്‍ ഒന്നിച്ചുകൂടുമ്പോള്‍ എന്റെ ഓര്‍മ്മയ്ക്കായി ചെയ്യുവിന്‍.}\\~\\
\Notice{എലാവരും കുനിഞ്ഞ് ആചാരം ചെയ്യുന്നു}
\end{frame}

\begin{frame}[allowframebreaks]
\Priest{
കര്‍ത്താവേ, നീ ഞങ്ങളോടു കല്പിച്ചതുപോലെ, എളിയവരും 
ബലഹീനരും ആകുലരുമായ ഈ ദാസര്‍ നിന്റെ തിരുസന്നിധിയില്‍ 
ഒരുമിച്ചു കൂടിയിരിക്കുന്നു. നന്ദി പ്രകാശിപ്പി ക്കുവാന്‍ കഴിയാത്ത വിധം അത്ര 
വലിയ അനുഗ്രഹമാണ് നീ ഞങ്ങള്‍ക്കു നല്കിയിരിക്കുന്നത്. നിന്റെ 
ദൈവിക ജീവനില്‍ ഞങ്ങളെ പങ്കുകാരാക്കുവാന്‍  \\ \framebreak 
നീ ഞങ്ങളുടെ മനുഷ്യസ്വഭാവം സ്വീകരിക്കുകയും അധ:പതിച്ചുപോയ ഞങ്ങളെ 
സമുദ്ധരിക്കുകയും മൃതരായ ഞങ്ങളെ ജീവിപ്പിക്കുകയും ചെയ്തു. 
പാപികളായ ഞങ്ങളെ കടങ്ങള്‍ ക്ഷമിച്ചു വിശുദ്ധീകരിച്ചു; 
ഞങ്ങളുടെ ബുദ്ധിക്കു പ്രകാശം നല്കി. ഞങ്ങളുടെ ശത്രുക്കളെ 
പരാജിതരാക്കി. ഞങ്ങളുടെ ബലഹീനമായ പ്രകൃതിയെ നിന്റെ \\ \framebreak
സമൃദ്ധമായ അനുഗ്രഹത്താല്‍ മഹത്വമണിയിക്കുകയും ചെയ്തു.

നീ ഞങ്ങള്‍ക്കു നല്കിയ എല്ലാ സഹായങ്ങള്‍ക്കും അനുഗ്രഹങ്ങള്‍ക്കുമായി 
ഞങ്ങള്‍ നിനക്ക് സ്തുതിയും ബഹുമാനവും കൃതജ്ഞതയും ആരാധനയും 
സമര്‍പ്പിക്കുന്നു. ഇപ്പോഴും എപ്പോഴും എന്നേക്കും.}
\Ammen
\end{frame}

\begin{frame}
\Server{ഹൃദയം കൊണ്ട് നിങ്ങള്‍ പ്രാര്‍ത്ഥിക്കുവിന്‍.\\ \Peace}
\end{frame}

\begin{frame}[allowframebreaks]
\frametitle{മാദ്ധ്യസ്ഥപ്രാര്‍ത്ഥന}
\Priest{
കര്‍ത്താവേ, ശക്തനായ ദൈവമേ, പ്രധാനാചാര്യനും സാര്‍വ്വത്രികസഭയുടെ 
തലവനും ഭരണാധികാരിയു മായ റോമിലെ മാര്‍ Francis പാപ്പയ്ക്കുവേണ്ടിയും 
ഞങ്ങളുടെ സഭയുടെ തലവനും പിതാവുമായ മാര്‍ ... 
ശ്രേഷ്ഠമേത്രാപ്പോലീത്തായ്ക്കുവേണ്ടിയും അതിരൂപാദ്ധ്യക്ഷന്‍ മാര്‍ ... 
മെത്രാപ്പോലീത്തക്കു വേണ്ടിയും ഞങ്ങളുടെ രൂപതാധ്യക്ഷനും പിതാവുമായ 
മാര്‍ ... മെത്രാനുവേണ്ടിയും വിശുദ്ധ കത്തോലിക്കാസഭ മുഴുവനുവേണ്ടിയും 
പുരോഹിതന്മാര്‍, ഭരണകര്‍ത്താക്കള്‍, മേലധികാരികള്‍ എന്നിവര്‍ക്കു 
വേണ്ടിയും ഈ കുര്‍ബാന സ്വീകരിക്കണമേ.}\par
\People{കര്‍ത്താവേ, ഈ കുര്‍ബാന സ്വീകരിക്കണമേ.}\\ \framebreak
\Priest{
എല്ലാ പ്രവാചകന്മാരുടെയും ശ്ലീഹന്‍ മാരുടെയും രക്തസാക്ഷികളുടെയും 
വന്ദകരുടെയും ബഹുമാനത്തിനു വേണ്ടിയും അങ്ങയുടെ സന്നിധിയില്‍ 
പ്രീതിജനകമായ വിധം വര്‍ത്തിച്ച നീതിമാന്മാരും വിശുദ്ധരുമായ എല്ലാ 
പിതാക്കന്മാര്‍ക്കു വേണ്ടിയും ഈ കുര്‍ബാന സ്വീകരിക്കണമേ.}\par
\People{കര്‍ത്താവേ, ഈ കുര്‍ബാന സ്വീകരിക്കണമേ.}\\ \framebreak
\Priest{
ക്ലേശിതരും ദുഖിതരും ദരിദ്രരും പീഡിതരും രോഗികളും ആകുലരുമായ 
എല്ലാവര്‍ക്കുവേണ്ടിയും ഞങ്ങളുടെ ഇടയില്‍ നിന്ന് അങ്ങയുടെ നാമത്തില്‍ 
വേര്‍പിരിഞ്ഞുപോയ എല്ലാ മരിച്ചവര്‍ക്കുവേണ്ടിയും  അങ്ങയുടെ കാരുണ്യത്തെ 
പ്രത്യാശാപൂര്‍വ്വം കാത്തിരിക്കുന്ന ഈ ജനത്തിനു  \\ \framebreak
വേണ്ടിയും അയോഗ്യനായ എനിക്കു വേണ്ടിയും ഈ കുര്‍ബാന സ്വീകരിക്കണമേ.}\par
\People{കര്‍ത്താവേ, ഈ കുര്‍ബാന സ്വീകരിക്കണമേ.}\\
\end{frame}

\begin{frame}[allowframebreaks]
\frametitle{നാലാം പ്രണാമജപം}
\Priest{
കര്‍ത്താവായ ദൈവമേ, അങ്ങ് ഞങ്ങളെ
......\\
സ്തുതിക്കുകയും അനുഷ്ടിക്കുകയും ചെയ്യുന്നു.}
\end{frame}

% maybe not required
\begin{frame}
\Server{നിങ്ങള്‍ നിശബ്ദരായി ആദരപൂര്‍വ്വം പ്രാര്‍ത്ഥിക്കുവിന്‍.}\\ \Peace \\~\\
\Priest{കര്‍ത്താവേ, നിന്റെ പരിശുദ്ധാത്മാവ് എഴുന്നള്ളി വരട്ടെ.}
\end{frame}

\begin{frame}
\Priest{
നിന്റെ ദാസരുടെ ഈ കുര്‍ബാനയില്‍ അവിടുന്ന് ആവസിച്ചു ഇതിനെ 
ആശീര്‍വ്വദിക്കുകയും പവിത്രീ കരിക്കുകയും ചെയ്യട്ടെ. ഇത് ഞങ്ങള്‍ക്ക് 
കടങ്ങളുടെ പൊറുതിക്കും പാപങ്ങളുടെ മോചനത്തിനും മരിച്ചവരുടെ 
ഉയിര്‍പ്പിലുള്ള വലിയ പ്രത്യാശയ്ക്കും നിന്നെ പ്രീതിപ്പെടുത്തിയ എല്ലാവരോടും 
ഒന്നിച്ച് സ്വര്‍ഗ്ഗരാജ്യത്തില്‍ നവമായ ജീവിതത്തിന്നും കാരണമാകട്ടെ.}
\end{frame}

\begin{frame}
\Priest{
കര്‍ത്താവായ ദൈവമേ, ഞങ്ങള്‍ക്കു വേണ്ടിയുള്ള അങ്ങയുടെ മഹനീയവും 
വിസ്മയാവഹവുമായ ഈ രക്ഷാ പദ്ധതിയെക്കുറിച്ച് അങ്ങയെ ഞങ്ങള്‍ നിരന്തരം 
പ്രകീര്‍ത്തിക്കുന്നു. അങ്ങയുടെ അഭിക്ഷിക്തന്റെ അമൂല്യരക്തത്താല്‍ രക്ഷിക്കപ്പെട്ട 
സഭയില്‍ സന്തോഷത്തോടും പ്രത്യാശയോടും കൂടെ ഞങ്ങള്‍ അങ്ങേക്കു 
കൃതജ്ഞത അര്‍പ്പിക്കുകയും ചെയ്യുന്നു.}
\end{frame}

\begin{frame}
\Priest{
സജീവവും പരിശുദ്ധവും ജീവദായകവുമായ അങ്ങയുടെ നാമത്തിന് 
സ്തുതിയും ബഹുമാനവും കൃതജ്ഞതയും ആരാധനയും ഞങ്ങള്‍ സമര്‍പ്പിക്കുന്നു.
ഇപ്പോഴും + എപ്പോഴും എന്നേക്കും.}\par \Ammen
\end{frame}

\begin{frame}
\frametitle{അനുരഞ്ജനശുശ്രൂഷ}
\Priest{ദൈവമേ, അങ്ങയുടെ കാരുണ്യത്തി നൊത്തവിധം എന്നോടു ദയ തോന്നണമേ.}\par
\People{അങ്ങയുടെ കാരുണ്യാതിരേകത്തി നനുസൃതമായി എന്റെ പാപങ്ങള്‍ മായിച്ചുകളയണമേ.}
\end{frame}

\begin{frame}
\Priest{
കര്‍ത്താവേ, അനുഗ്രഹിക്കണമേ. ഞങ്ങളുടെ കര്‍ത്താവായ ദൈവമേ, 
ഞങ്ങള്‍ അയോഗ്യരാകുന്നു. ഞങ്ങള്‍ തീര്‍ത്തും അയോഗ്യരാകുന്നു.
എങ്കിലും സ്തുത്യര്‍ഹവും പരിശുദ്ധവും ജീവദായകവും ദൈവികവുമായ
ഈ രഹസ്യങ്ങളിലേക്ക് അങ്ങയുടെ കാരുണ്യം ഞങ്ങളെ അടുപ്പിക്കുന്നു.}
\end{frame}

\begin{frame}[allowframebreaks]
\frametitle{വിഭജനശുശ്രൂഷ}
\Priest{
ഞങ്ങളുടെ കര്‍ത്താവീശോമിശിഹായെ, നിന്റെ തിരുനാമാത്തിനു സ്തുതിയും നാഥനായ 
നിനക്ക് ആരാധനയും എപ്പോഴും ഉണ്ടായിരിക്കട്ടെ. സജീവവും ജീവദായകവുമായ
ഈ അപ്പം സ്വര്‍ഗ്ഗത്തില്‍ നിന്ന് ഇറങ്ങിയതും ലോകത്തിനു മുഴുവന്‍ ജീവന്‍
നല്‍കുന്നതുമാകുന്നു. ഇതു ഭക്ഷിക്കുന്നവര്‍ മരിക്കുകയില്ല. പ്രത്യുത, പാപമോചനവും
രക്ഷയും പ്രാപിക്കുകയും നിത്യം ജീവിക്കുകയും ചെയ്യും.}
\end{frame}

\begin{frame}[allowframebreaks]
% \frametitle{ആരാധനാഗീതം} % <- no space to add
\Server{ഞാന്‍ സ്വര്‍ഗ്ഗത്തില്‍ നിന്നിറങ്ങിയ ജീവനുള്ള അപ്പമാകുന്നു.}\par
\People{
രക്ഷകനീശോതന്‍, ശിഷ്യരെയറിയിച്ച
ദിവ്യരഹസ്യമിതാ:
സ്വര്‍ഗ്ഗത്തില്‍നിന്നാഗതമാം 
ജീവന്‍ നല്കിടുമപ്പം ഞാന്‍,
സ്നേഹമോടെന്നെ കൈക്കൊള്‍വോനെന്നില്‍ 
നിത്യം ജീവിക്കും,
നേടുമവന്‍ സ്വര്‍ഗ്ഗം നിശ്ചയമായ്.}
\end{frame}

\begin{frame}
\Priest{
നമ്മുടെ കര്‍ത്താവീശോമിശിഹായുടെ കൃപയും പിതാവായ ദൈവത്തിന്റെ 
സ്നേഹവും പരിശുദ്ധാത്മാവിന്റെ സഹവാസവും നാമെല്ലാവരോടും കൂടെ ഉണ്ടായിരിക്കട്ടെ.
ഇപ്പോഴും + എപ്പോഴും എന്നേക്കും.}\par \Ammen
\end{frame}

\begin{frame}[allowframebreaks]
\frametitle{അനുരഞ്ജനഗീതം}
\Server{
നമ്മുടെ രക്ഷകന്റെ അമൂല്യമായ ശരീരരക്തങ്ങളെ നമുക്കെല്ലാവര്‍ക്കും 
ഭക്ത്യാദരങ്ങളോടെ സമീപിക്കാം. അനുതാപത്തില്‍ നിന്നുളവാകുന്ന
ശരണത്തോടെ അപരാധങ്ങളില്‍ നിന്നു പിന്തിരിഞ്ഞുകൊണ്ടും 
പാപങ്ങളെ ക്കുറിച്ച് പശ്ചാത്തപിച്ചുകൊണ്ടും സഹോദരങ്ങളുടെ 
തെറ്റുകള്‍ ക്ഷമിച്ചുകൊണ്ടും നമുക്ക് സകലത്തിന്റെയും നാഥനായ 
ദൈവത്തോട് കൃപയും പാപമോചനവും യാചിക്കുകയും ചെയ്യാം.}
\end{frame}

\begin{frame}[allowframebreaks]
\ForgiveUs\par
\Server{ഭിന്നതകളും കലഹങ്ങളും വെടിഞ്ഞ് നമ്മുടെ മനസ്സാക്ഷിയെ ശുദ്ധീകരിക്കാം.}\\
\People{കര്‍ത്താവേ, അങ്ങയുടെ ... }\par
\Server{ശത്രുതയിലും വിദ്വേഷത്തില്ലുംനിന്ന് നമ്മുടെ ആത്മാക്കളെ വിമുക്തമാക്കാം.}\\
\People{കര്‍ത്താവേ, അങ്ങയുടെ ... }\\ \framebreak
\Server{വിശുദ്ധ കുര്‍ബാന സ്വീകരിക്കുകയും പരിശുദ്ധാത്മാവിനാല്‍ പവിത്രീകരിക്കപ്പെടുകയും ചെയ്യാം.}\\
\ForgiveUs\par
\Server{യോജിപ്പോടും ഐക്യത്തോടുംകൂടെ ദിവ്യരഹസ്യങ്ങളില്‍ പങ്കുകൊള്ളാം.}\\
\People{കര്‍ത്താവേ, അങ്ങയുടെ ... }\\ \framebreak
\Server{കര്‍ത്താവേ, ഇവ ഞങ്ങളുടെ ശരീരങ്ങളുടെ ഉയിര്‍പ്പിനും ആത്മാക്കളുടെ രക്ഷയ്ക്കും കാരണമാകട്ടെ.}\par
\People{നിത്യജീവനും കാരണമാകട്ടെ, എന്നേക്കും} \Ammen
\end{frame}

% OR
\iftrue
\begin{frame}[allowframebreaks]
മനസ്സില്‍ നിറയും മലിനതയെല്ലാം
കഴുകി ധന്യരാകാം.

അപരാധങ്ങള്‍ നീക്കണമേ.
പാപകടങ്ങള്‍ മായ്ക്കണമേ.

ശത്രുതയെല്ലാം നീക്കി നിതാന്തം
ഹൃദയം ദീപ്തമാക്കാം.

അപരാധങ്ങള്‍ നീക്കണമേ.
പാപകടങ്ങള്‍ മായ്ക്കണമേ.

ഹൃദയമൊരുക്കാം നാഥനു നിത്യം
നിറയാം അത്മാവാല്‍.

അപരാധങ്ങള്‍ നീക്കണമേ.
പാപകടങ്ങള്‍ മായ്ക്കണമേ.

ഏകമനസ്സായ് ദിവ്യരഹസ്യം
നിതരാമുള്‍ക്കൊള്ളാം.

അപരാധങ്ങള്‍ നീക്കണമേ.
പാപകടങ്ങള്‍ മായ്ക്കണമേ.

ആത്മാവിന്‍ പരിരക്ഷണമതിനീ
നിണവും തിരുമെയ്യും

അപരാധങ്ങള്‍ നീക്കണമേ
പാപകടങ്ങള്‍ മായ്ക്കണമേ

ഉത്ഥാനത്തിനുമതുപോല്‍ ജീവനും
കാരണമായിടുവാന്‍

അപരാധങ്ങള്‍ നീക്കണമേ
പാപകടങ്ങള്‍ മായ്ക്കണമേ
\end{frame}
\fi

\begin{frame}
\LetsPray
\end{frame}

\begin{frame}[allowframebreaks]
\frametitle{ദൈവൈക്യശുശ്രൂഷ}
\Priest{
കര്‍ത്താവേ, ഞങ്ങളുടെ ഇടയില്‍ നിന്റെ സമാധാനവും ശാന്തിയും പുലര്‍ത്തണമേ.
ഞങ്ങളുടെ നാവുകള്‍ നിന്റെ സത്യം പ്രഘോഷിക്കട്ടെ. ഞങ്ങളുടെ വദനങ്ങളെ ഞങ്ങള്‍ 
പുതിയ വീണകളാക്കുകയും ജ്വലിക്കുന്ന അധരങ്ങള്‍കൊണ്ട് ഞങ്ങള്‍ സ്തോത്രം 
ആലപിക്കുകയും ചെയ്യുമ്പോള്‍ നിന്റെ കുരിശു ഞങ്ങക്കു സംരക്ഷണമായി ഭവിക്കട്ടെ. 
നിങ്ങള്‍ പ്രാര്‍ത്ഥിക്കുമ്പോള്‍ ഇപ്രകാരം പ്രാര്‍ത്ഥിക്കുവിന്‍ എന്നുപദേശിച്ചുകൊണ്ട് 
നീ ശിഷ്യന്മാരെ പഠിപ്പിച്ച വിശിഷ്ട്ടമായ പ്രാര്‍ത്ഥന നിന്നില്‍നിന്നു ലഭിച്ച 
മനോവിശ്വാസത്തോടെ ആവര്‍ത്തിക്കുവാന്‍ ഞങ്ങളെ യോഗ്യരാക്കണമേ.}
\end{frame}

\begin{frame}[allowframebreaks]
\frametitle{കര്‍ത്ത പ്രാര്‍തന}
\Priest{സ്വര്‍ഗ്ഗസ്ഥനായ ഞങ്ങളുടെ പിതാവേ,}\\
\People{
അങ്ങയുടെ നാമം പൂജിതമാകണമേ./ അങ്ങയുടെ രാജ്യം
വരണമേ. അങ്ങയുടെ തിരുമനസ്സ് സ്വര്‍ഗ്ഗത്തിലെപ്പോലെ 
ഭൂമിയിലുമാകണമേ./ ഞങ്ങള്‍ക്ക് ആവശ്യകമായ ആഹാരം 
ഇന്ന് ഞങ്ങള്‍ക്കു തരണമേ./ ഞങ്ങളുടെ കടക്കരോട് 
ഞങ്ങള്‍ ക്ഷമിച്ചിരിക്കുന്നത് പോലെ/ ഞങ്ങളുടെ കടങ്ങളും 
പാപങ്ങളും/ ഞങ്ങളോടും ക്ഷമിക്കണമേ./ ഞങ്ങളെ 
പ്രലോഭനത്തില്‍ ഉള്‍പ്പെടുത്തരുതേ. ദുഷ്ട്ടാരൂപിയില്‍ 
നിന്ന് ഞങ്ങളെ രക്ഷിക്കണമേ./ എന്തുകൊണ്ടെന്നാല്‍, 
രാജ്യവും ശക്തിയും മഹത്ത്വവും എന്നേക്കും അങ്ങയുടേതാകുന്നു.}
\Ammen
\end{frame}

\begin{frame}[allowframebreaks]
\Priest{
കര്‍ത്താവേ, ശക്തനായ സര്‍വ്വേശ്വരാ, നല്ലവനായ ദൈവമേ,
കൃപാപൂര്‍ണ്ണനായ ഞങ്ങളുടെ പിതാവേ, അങ്ങയുടെ കാരുണ്യം 
ഞങ്ങള്‍ കേണപേക്ഷിക്കുന്നു. ഞങ്ങളെ പ്രലോഭനത്തില്‍ 
ഉള്‍പ്പെടുത്തരുതേ. ദുഷ്ടാരൂപിയില്‍നിന്നും അവന്റെ 
സൈന്യങ്ങളില്‍നിന്നും ഞങ്ങളെ രക്ഷിക്കണമേ. \framebreak
എന്തുകൊണ്ടെന്നാല്‍, സ്വര്‍ഗ്ഗത്തിലും ഭൂമിയിലും രാജ്യവും 
ശക്തിയും പ്രാബല്യവും അധികാരവും അങ്ങയുടേതാകുന്നു.
ഇപ്പോഴും + എപ്പോഴും എന്നേക്കും.}\\\Ammen
\end{frame}

\begin{frame}
\frametitle{സമാധാനാശംസ}
\PPeace\\~\\
\People{അങ്ങയോടും അങ്ങയുടെ ആത്മാവോടുംകൂടെ.}
\end{frame}

\begin{frame}
\Priest{വിശുദ്ധ കുര്‍ബാന വിശുദ്ധ ജനത്തിനുള്ളതാകുന്നു.}\par
\People{
ഏക പിതാവു പരിശുദ്ധനാകുന്നു./ ഏകപുത്രന്‍ പരിശുദ്ധനാകുന്നു./
ഏകറൂഹാ പരിശുദ്ധനാകുന്നു./
പിതാവിനും പുത്രനും പരിശുദ്ധാത്മാവിനും എന്നേക്കും സ്തുതി.}\\\Ammen
\end{frame}


% TODO:
% According to kalam

\begin{frame}[allowframebreaks]
\frametitle{ദിവ്യകാരുണ്യസ്തുതി-നോബുകാലം}
\People{ഇതാ ജീവന്റെ ദിവ്യൗഷധം നല്‍കപ്പെടുന്നു. സഹോദരേ, നിങ്ങള്‍ ഇതു സ്വീകരിച്ച് ജീവന്‍ പ്രാപിക്കൂവിന്‍.}
\end{frame}


\begin{frame}
\frametitle{ദിവ്യകാരുണ്യസ്തുതി}
\DivyaKarunyaSthuthi\par
\People{
സഭയിൽ അവിടുത്തേക്കു സ്തുതിയുണ്ടായിരിക്കട്ടെ.
എല്ലാ നിമിഷവും എല്ലാ സമയവും / അവിടുത്തെ കൃപയും കാരുണ്യവും /
നമ്മുടെമേൽ ഉണ്ടായിരിക്കട്ടെ.}
\end{frame}

\iffalse
\begin{frame}
\frametitle{ദിവ്യകാരുണ്യസ്തുതി}
\DivyaKarunyaSthuthi\par

% kalam ?
\People{
സ്വര്‍ഗ്ഗത്തില്‍നിന്നൊരു സ്വരം മാറ്റൊലിക്കൊണ്ടു: “ഇവന്‍ എന്‍റെ പ്രിയപുത്രനാകുന്നു.
ഇവനില്‍ ഞാന്‍ സംപ്രീതനായിരി ക്കുന്നു.” ലോകരക്ഷനെ നമുക്കു ഒന്നുചേര്‍ന്ന് സ്തുതിച്ചാരാധിക്കം.}

% kalam ?
\People{
കര്‍ത്താവേ, നിന്‍റെ തിരുശരീരരക്തങ്ങള്‍ സ്വീകരിക്കുവാന്‍ ഞങ്ങളെ നീ അനുഗ്രഹിച്ചു.
നിന്റെ രാജ്യത്തിലേക്കു കാരുണ്യപൂര്‍വ്വം ഞങ്ങളെ ക്ഷണിക്കുകയും ചെയ്തു.}

% kalam ?
\People{
കര്‍ത്താവായ ദൈവം മനുഷ്യസ്വഭാവം സ്വീകരിച്ച് നമ്മിലൊരുവനായി.
തന്റെ ശരീരവും രക്തവും ദാനം ചെയ്തുകൊണ്ട് അവിടുന്ന് നമ്മെ തന്റെ 
മഹത്വത്തില്‍ പങ്കുകാരാക്കി. ഈ മഹാദാനം സ്നേഹപൂര്‍വ്വം സ്വീകരിച്ചുകൊണ്ട് 
നല്ലവനായ കര്‍ത്താവിനെ നമുക്കു സ്തുതിക്കാം.}

% kalam ?
\People{
നമ്മെ വിശുദ്ധീകരിക്കുന്ന അമൂല്യമായ ശരീരരക്തങ്ങള്‍ നല്കപ്പെടുന്നു. 
സകലത്തിന്റെയും നാഥാ, നീ പരിശുദ്ധന്‍, പരിശുദ്ധന്‍, പരിശുദ്ധന്‍, 
എന്നു നമുക്കു പാടി സ്തുതിക്കാം. പാപമോചനത്തിനായി ഇതാ 
സകലജനതകള്‍ക്കും സ്വര്‍ഗ്ഗീയ അപ്പം നല്കപ്പെടുന്നു. നമ്മുക്ക് ഈ 
ദിവ്യരഹസ്യങ്ങള്‍ സ്വീകരിച്ച് മാലാഖമാരോടൊത്ത് അവിടുത്തെ ഹല്ലേല്ലുയ്യ പടി സ്തുതിക്കാം.}

% kalam ?
\People{
വിശ്വാസമാകുന്ന പാറമേല്‍ തന്റെ സഭയെ പണിതുയര്‍ത്തുകയും ശ്ലീഹന്മാരെയും 
പ്രവാചകന്മാരെയും സത്യ പ്രബോധകരെയും സഭയ്ക്കു നല്കുകയും ചെയ്യുന്നവന്‍ വാഴ്ത്തപ്പെട്ടവനാകുന്നു.}

% നോബുകാലം
\People{
ഇതാ ജീവന്റെ ദിവ്യൗഷധം നല്‍കപ്പെടുന്നു. സഹോദരേ, നിങ്ങള്‍ ഇതു സവീകരിച്ച് ജീവന്‍ പ്രാപിക്കൂവിന്‍.}
\end{frame}
\fi

\begin{frame}
\frametitle{പരിശുദ്ധ കുര്‍ബാന സ്വീകരണം} 
\Server{
നമ്മെ ജീവിപ്പിക്കുന്ന കര്‍ത്താവീശോമിശിഹായുടെ കൃപാവരം 
അവിടുത്തെ കാരുണ്യത്താല്‍ നാമെല്ലാവരില്ലും സമ്പൂര്‍ണമാകട്ടെ.}\par
\People{എപ്പോഴും എന്നേക്കും} \Ammen
\end{frame}

% TODO:
% According to kalam

\begin{frame}[allowframebreaks]
\frametitle{ദിവ്യകാരുണ്യ സ്വീകരണഗീതം\\നോബുകാലം}
\Server{വരുവി, ഉപവാസവും പ്രായശ്ചിത്തവും വഴി പാപക്കറകള്‍ 
കഴുകിക്കളയുന്ന ആത്മീയശക്തി നമുക്ക് പ്രാപിക്കാം. ഭയഭക്തികളോടും 
സ്നേഹത്തോടും കൂടി നമുക്കവിടുത്തെ തിരുശരീരരക്തങ്ങള്‍ സ്വീകരിക്കാം.} 
\end{frame}


\begin{frame}
\frametitle{ദിവ്യകാരുണ്യ സ്വീകരണഗീതം}
\Server{
സഹോദരരേ, സ്വര്‍ഗ്ഗരാജ്യത്തിലേക്കുള്ള വിശ്വാസത്തോടെ ദൈവപുത്രന്റെ ശരീരം 
സ്വീകരിക്കുവാനും അവിടുത്തെ രക്തം പാനം ചെയ്യുവാനും തിരുസ്സഭ നിങ്ങളെ ക്ഷണിക്കുന്നു.}
\end{frame}

% OR

\begin{frame}
\frametitle{ദിവ്യകാരുണ്യ സ്വീകരണഗീതം}
\Server{
ജലത്താലും അരൂപിയാലും നാം മിശിഹായുടെ മൌതികശരീരത്തിന്റെ 
അംഗങ്ങളായിരിക്കുന്നു. അവിടുത്തെ ശരീരവും രക്തവും നമ്മെ പരിപോഷിപ്പിക്കട്ടെ.}
\end{frame}

\begin{frame}[allowframebreaks]
\frametitle{കൃതജ്ഞതാപ്രകാശനം}
\Server{
പരിശുദ്ധാത്മാവിന്റെ കൃപാവരത്താല്‍ സ്തുത്യര്‍ഹവും പരിശുദ്ധവും 
ജീവദായകവും ദൈവികവുമായ ഈ രഹസ്യങ്ങളെ സമീപിച്ച്
ഇവയില്‍ പങ്കുകൊള്ളുവാന്‍ യോഗ്യരാക്കപ്പെട്ട നമുക്കെല്ലാവര്‍ക്കും 
ഇവയുടെ ദാതാവായ ദൈവത്തിന് സ്തുതിയും കൃതജ്ഞതയും സമര്‍പ്പിക്കാം.}\\ \framebreak
\People{അവര്‍ണ്ണനീയമായ ഈ ദാനത്തെക്കുറിച്ച് കര്‍ത്താവേ, അങ്ങേക്കു സ്തുതി.}\\~\\
\LetsPray
\end{frame}

% Sundays
\begin{frame}[allowframebreaks]
\Priest{
കര്‍ത്താവായ ദൈവമേ, അങ്ങയുടെ മഹനീയമായ നാമത്തിനു സ്തുതിയും 
ആരാധനയും കൃതജ്ഞതയും സമര്‍പ്പിക്കുക എല്ലാ നാഴികകളിലും 
എല്ലാക്കാലങ്ങളിലും യുക്തവും ആകുന്നു. അങ്ങയുടെ ജീവദായകവും 
ദൈവികവും ആയ വചനങ്ങളുടെ മാധുര്യം ആസ്വദിക്കുവാനും മാലാഖമാരോടുകൂടെ 
തിരുനാമത്തെ സ്തുതിക്കുവാനും അങ്ങയുടെ ദാനമായ ദിവ്യരഹസ്യ ങ്ങളില്‍ 
പങ്കുകൊള്ളുവാനും മഹോന്നതമായ ദൈവമേ, അങ്ങേക്കു സ്തുതിയുടെയും 
കൃതജ്ഞതയുടെയും കീര്‍ത്തനങ്ങള്‍ ഇടവിടാതെ ആലപിക്കുവാനും 
ബലഹീനരായ ഞങ്ങളെ കാരുണ്യപൂര്‍വ്വം അങ്ങു യോഗ്യരാക്കി. 
\SignOfCross}
\Ammen\\ \People{ കര്‍ത്താവേ ആശീര്‍വ്വദിക്കണമേ.}\\
\end{frame}

\begin{frame}[allowframebreaks]
\Priest{
സകലത്തെയും പവിത്രീകരിക്കുന്ന തന്റെ അമൂല്യമായ ശരീരരക്തങ്ങള്‍ 
സ്വീകരിക്കുവാന്‍ നമ്മുടെ കര്‍ത്താവും ദൈവവും രാജാവും രക്ഷകനും 
ജീവദാതാവുമായ മിശിഹാ കൃപാപൂര്‍വ്വം നമ്മെ യോഗ്യരാക്കി.
വിചാരങ്ങളും വചനങ്ങളും പ്രവര്‍ത്തികളും വഴി തന്നെ പ്രസാദിപ്പിക്കുവാന്‍ 
അവിടുന്ന് നമ്മെ അനുഗ്രഹിക്കട്ടെ. കര്‍ത്താവേ, ഞങ്ങള്‍
സ്വീകരിച്ചതും സ്വീകരിക്കുന്നതുമായ ഈ അച്ചാരം നിന്റെ 
കാരുണ്യാതി രേകത്താല്‍ ഞങ്ങളുടെ കടങ്ങളുടെ പോറുതിക്കും 
പാപങ്ങളുടെ മോചന ത്തിനും മരിച്ചവരുടെ ഉയര്‍പ്പിലുള്ള 
പ്രത്യാശയ്ക്കും നിന്റെ സന്നിധിയില്‍ പ്രീതിജനകമായവിധം വര്‍ത്തിച്ച 
എല്ലാവരോടുമോന്നിച്ച് സ്വര്‍ഗ്ഗ രാജ്യത്തില്‍ നവമായ ജീവിതത്തിന്നും 
കാരണമാകട്ടെ. സകലത്തിന്റെയും നാഥാ, എന്നേക്കും,}
\Ammen\\
\end{frame}

% TODO: Other days

\begin{frame}[allowframebreaks]
\frametitle{സമാപനാശീര്‍വ്വാദം}
\Priest{
കര്‍ത്താവാം മിശിഹാവഴിയായ്\\
ദിവ്യാത്മാവിന്‍ ദാനങ്ങള്‍ \\
സ്നേഹപിതാവാം സകലേശന്‍\\
വിരവൊടു നമ്മില്‍ വര്‍ഷിച്ചു.

സ്വര്‍ല്ലോകത്തിന്‍  മഹിമയ്ക്കായ്‌ \\
ദൈവം നമ്മെ വിളിച്ചല്ലോ \\
അക്ഷയസൗഭഗമാര്‍ന്നിടുവാന്‍\\
ദൈവം നമ്മെ നയിച്ചല്ലോ.}
\Ammen\\ \framebreak

\Priest{
എന്റെ ശരീരം ഭക്ഷിക്കും\\
എന്‍ രക്തം പാനം ചെയ്യും\\
മാനവനെന്നില്‍ നിവസിക്കും\\
അവനില്‍ ഞാനും, നിശ്ചയമായ്

വിധിയില്‍ വീഴാതവനെ ഞാന്‍\\
അന്തിമദിവസമുയര്‍പ്പിക്കും\\
നിത്യായുസ്സവനേകും ഞാന്‍ \\
ഏവം നാഥനരുള്‍ചെയ്തു.}
\Ammen\\ \framebreak

\Priest{
അരുളിയപോലിന്നീ ബലിയില്‍\\
പങ്കാളികളാമഖിലര്‍ക്കും\\
ദൈവം കനിവാര്‍ന്നരുളട്ടെ\\
ദിവ്യാനുഗ്രഹമെന്നെന്നും.

നമ്മുടെ ജീവിതപാതകളില്‍ \\
ആവശ്യകമാം ദാനങ്ങള്‍\\
നല്കിയനുഗ്രഹമരുളട്ടെ\\
ദൈവം കരുണയോടെന്നെന്നും.}
\Ammen\\ \framebreak

\Priest{
ജീവന്‍ നല്കും ദൈവീകമാം \\
ശുദ്ധിയെഴുന്ന രഹസ്യങ്ങള്‍ \\
കൈക്കൊണ്ടവരാമഖിലരിലും\\
ദൈവം വരനിര ചൊരിയട്ടെ.

കുരിശടയാളം വഴിയായ് നാം\\
സംരക്ഷിതരായിത്തീരട്ടെ\\
മുദ്രിതരായി ഭവിക്കട്ടെ +\\
ഇപ്പോഴുമെപ്പൊഴുമെന്നേക്കും.}
\Ammen\\ \framebreak
\end{frame}

% TODO Other days

\begin{frame}
\frametitle{വിടവാങ്ങല്‍ പ്രാര്‍ത്ഥന}
\Priest{വിശുദ്ധീകരണത്തിന്‍റെ ബലിപീഠമേ സ്വസ്തി!}
\People{
ഞങ്ങളുടെ കര്‍ത്താവിന്റെ കബറിടമേ, സ്വസ്തി!/ നിന്നില്‍നിന്നു 
ഞാന്‍ സ്വീകരിച്ച ഈ കുര്‍ബാന/ കടങ്ങളുടെ പൊറുതിക്കും/ 
പാപങ്ങളുടെ മോചനത്തിനും കാരണമാകട്ടെ./ ഇനി ഒരു 
ബലി അര്‍പ്പിക്കുവാന്‍/ ഞാന്‍ വരുമോ ഇല്ലയോ/ എന്നറിഞ്ഞുകൂടാ.}
\end{frame}

\begin{frame}[allowframebreaks]
\frametitle{വൈദികര്‍ക്കുവേണ്ടിയുള്ള പ്രാര്‍ത്ഥന}
\People{
നിത്യപുരോഹിതനായ ഈശോ, / അങ്ങേ ദാസന്മാരായ
വൈദികര്‍ക്ക്/ യാതൊരാപത്തും വരാതെ/ അങ്ങേ തിരുഹൃദയത്തില്‍ 
അഭയം നല്കണമേ. അങ്ങയുടെ പരിശുദ്ധമായ ശരീരത്തെ/ ദിവസംതോറും
എടുക്കുന്ന അവരുടെ പൂശപ്പെട്ട കരങ്ങളെ/ മലിനമാക്കാതെ കാക്കണമേ. 
അങ്ങേ വിലയേറിയ \\ \framebreak തിരുരക്തത്താല്‍ നനയുന്ന 
അവരുടെ നാവുകളെ/ നിര്‍മ്മലമായി കാത്തുക്കൊള്ളണമേ. ശ്രേഷ്ഠമായ 
അങ്ങേ  പൌരോഹിത്യത്തിന്റെ/ മഹനീയമുദ്ര പതിച്ചിരിക്കുന്ന അവരുടെ 
ഹൃദയങ്ങളെ/ ലോകവസ്തുക്കളില്‍നിന്ന് അകറ്റുകയും/ വിശുദ്ധമായി 
കാത്തുകൊള്ളുകയും ചെയ്യണമേ. അങ്ങേ ദിവ്യസ്നേഹം/ അവരെ \\ \framebreak 
ലോകതന്ത്രങ്ങളില്‍നിന്നു സംരക്ഷിക്കട്ടെ. അവരുടെ പ്രയത്നങ്ങള്‍ 
ഫലസമൃദ്ധങ്ങളായി ഭവിക്കട്ടെ. അവരുടെ ശുശ്രൂഷ ലഭിക്കുന്നവര്‍/ 
ഇഹത്തില്‍ അവരുടെ ആനന്ദവും ആശ്വാസവും / പരത്തില്‍ 
നിത്യസൗഭാഗ്യത്തിന്റെ മകുടവും ആയിത്തീരട്ടെ. ആമ്മേന്‍. \\ \framebreak
ലോകരക്ഷകനായ ഈശോ, അങ്ങേ പുരോഹിതരേയും 
വൈദികശുശ്രൂഷകരെയും ശുദ്ധീകരിക്കണമേ.\\~\\
വൈദികരുടെ രാജ്ഞിയായ മറിയമേ,
വൈദികര്‍ക്കു വേണ്ടി പ്രാര്‍ത്ഥിക്കണമേ.}
\end{frame}

% Today Only
\iffalse
\begin{frame}[allowframebreaks]
\frametitle{ലുത്തിനിയ}
\People{
കര്‍ത്താവേ കനിയണമേ\\
മിശിഹായേ കനിയണമേ\\
കര്‍ത്താവേ ഞങ്ങളണയ്ക്കും\\
പ്രാര്‍ത്ഥന സദയം കേള്‍ക്കണമേ.

സ്വര്‍ഗ്ഗപിതാവാം സകലേശാ\\
ദിവ്യാനുഗ്രഹമേകണമേ\\
നരരക്ഷകനാം മിശിഹായേ\\
ദിവ്യാനുഗ്രഹമേകണമേ. \\ \framebreak

ദൈവാത്മാവാം സകലേശാ\\
ദിവ്യാനുഗ്രഹമേകണമേ\\
പരിപാവനമാം ത്രീത്വമേ\\
ദിവ്യാനുഗ്രഹമേകണമേ.

വിശുദ്ധിവിരിയും പാതകളിൽ\\
ചരിച്ചു നാഥനു സമ്മോദം\\
നിറച്ച പാവന ധീരാത്മാ,\\
പ്രാർത്ഥിക്കണമേ ഞങ്ങൾക്കായ്. \\ \framebreak

ക്രൂശിതനായൊരു ദൈവത്തിൻ\\
പൂജിതമാം തിരുമാതൃകയാൽ\\
കുരിശിൻ പാത പുണർന്നവനേ,\\
പ്രാർത്ഥിക്കണമേ ഞങ്ങൾക്കായ്.

ആനന്ദത്തിലുമല്ലലിലും\\
പ്രീതി മഹേശനു ചേർത്തൊടുവിൽ\\
നിതയ്കിരീടമണിഞ്ഞവേന,\\
പ്രാർത്ഥിക്കണമേ ഞങ്ങൾക്കായ്. \\ \framebreak

വിനീതനായ വിശുദ്ധാത്മാ,\\
വിരിഞ്ഞുെപാന്തിയ വിണ്മലേര,\\
നിറഞ്ഞനീതി പുണർന്നിടുവാൻ\\
പ്രാർത്ഥിക്കണമേ ഞങ്ങൾക്കായ്.

ലോകവുമതിനുെട ആശകളും\\
ധീരതയാർന്നു വെടിഞ്ഞവേന,\\
വിനയത്തിൻ തിരുമാതൃകേയ,\\
പ്രാർത്ഥിക്കണമേ ഞങ്ങൾക്കായ്. \\ \framebreak

ഉലകിൻ മോഹനദീപം നീ,\\
ധരയുടെ ശാശ്വതലവണം നീ,\\
മലയിലുയർന്നൊരു നഗരം നീ,\\
പ്രാർത്ഥിക്കണമേ ഞങ്ങൾക്കായ്.

ഉപവാസത്തിൻ പാതകളിൽ\\
പ്രാർത്ഥന തീർത്തൊരു ദീപവുമായ്\\
പാഞ്ഞു കിരീടമണിഞ്ഞവേന,\\
പ്രാർത്ഥിക്കണമേ ഞങ്ങൾക്കായ്. \\ \framebreak

സ്വർഗ്ഗതലത്തു മഹോന്നതമാം\\
നിക്ഷേപങ്ങൾ നിറച്ചവേന,\\
ചിന്തകൾ വിണ്ണിലുയർത്തിടുവാൻ\\
പ്രാർത്ഥിക്കണമേ ഞങ്ങൾക്കായ്.}
\end{frame}
\fi
\end{document}
